\section{Step 7: Admit a bounded record window}
The routine must decide whether the byte count is plausible before reading a
header. A record needs at least sixteen bytes and may contain at most 28\,672
payload bytes, so admit counts from 16 through 28\,688 inclusive. Treat this as an
unsigned 64-bit count. Narrowing it to 32 bits would make some very large counts
appear small. Check the image's low two address bits and the context's low three
bits without dereferencing them. The record address needs no alignment test.

Fill \texttt{record\_admit}. Keep the initialized failure return for rejected
arguments and branch to the supplied early return. An admitted call falls through
to the supplied 160-byte stack frame, which saves the return address and saved
registers before giving the four arguments persistent names in
\texttt{\$s0} through \texttt{\$s3}. A bad count must return even if the context
address is unusable.

\paragraph{Observe} Prepare \texttt{s07-1} and stop before
\texttt{record\_admitted}. The worked call has 27 record bytes. At \texttt{record\_admitted},
\texttt{\$a1} is 27 and the final alignment mask in \texttt{\$t0} is zero.
No frame or context bytes have changed.

\Needspace{7\baselineskip}
\paragraph{Check} Why is the record address allowed to be unaligned?
\begin{itemize}
  \item[A.] Its fields can be read as individual bytes.
  \item[B.] All records begin at a multiple of sixteen.
  \item[C.] A delta has no header.
\end{itemize}
